% 张三慧《学习辅导与习题解答》扫描 S 卷1第6章“思考题选答”原编号。
% 仅辅导书来源已视觉核验，V1主教材原题逐图交叉核验仍待；全部exam_ready=false。
\section*{第6章\quad 振动：辅导书已列思考题独立解析}
\paragraph{思考题 6.7}
自由振子满足 \(m\ddot x+kx=0\)，本征角频率 \(\omega_0=\sqrt{k/m}\)，且无外部持续供能；理想情形机械能守恒。稳定受迫振子通常满足 \(m\ddot x+b\dot x+kx=F_0\cos\Omega t\)，过渡过程衰减后以\emph{外力频率}\(\Omega\) 振动，振幅 \(A=F_0/\sqrt{(k-m\Omega^2)^2+b^2\Omega^2}\)，驱动力每周期提供的平均能量平衡阻尼耗散。即使二者形式上都为正弦，频率来源、能量维持及阻尼机制不同。

\paragraph{思考题 6.8}
低频、弹簧沿长度近似线性运动时，均匀弹簧质量 \(m_s\) 的动能可写成 \(\frac12(m_s/3)\dot x^2\)，等效附加惯性为 \(m_s/3\) 而非整个质量。故 \(\boxed{T\simeq2\pi\sqrt{(m+m_s/3)/k}}>2\pi\sqrt{m/k}\)。当 \(m_s\) 与重物质量同量级或发生高阶弹簧振动时，线性位移假设失效，应解分布质量弹性连续体的本征模方程；不能把教材“整个弹簧质量都加到球上”当成精确公式。
